\documentclass[12pt,a4paper]{article}

%% ── Packages ────────────────────────────────────────────────────────────────
\usepackage[utf8]{inputenc}
\usepackage[T1]{fontenc}
\usepackage{lmodern}
\usepackage{microtype}
\usepackage[margin=2.5cm]{geometry}
\usepackage{titlesec}
\usepackage{setspace}
\usepackage{amsmath,amssymb,amsthm}
\usepackage{mathtools}
\usepackage{bm}
\usepackage{listings}
\usepackage{xcolor}
\usepackage{booktabs}
\usepackage{longtable}
\usepackage{array}
\usepackage{hyperref}
\usepackage{fancyhdr}
\usepackage{enumitem}
\usepackage{tcolorbox}
\usepackage{mdframed}
\usepackage{algorithm}
\usepackage{algpseudocode}

%% ── Colours ─────────────────────────────────────────────────────────────────
\definecolor{multBlue}{RGB}{0,60,120}
\definecolor{multGold}{RGB}{180,140,0}
\definecolor{codeGray}{RGB}{245,245,250}
\definecolor{codeGreen}{RGB}{0,120,60}
\definecolor{codeRed}{RGB}{180,0,0}
\definecolor{codeBlue}{RGB}{0,80,180}
\definecolor{codePurple}{RGB}{120,0,160}

%% ── Hyperref ────────────────────────────────────────────────────────────────
\hypersetup{
  colorlinks = true,
  linkcolor  = multBlue,
  urlcolor   = multBlue,
  citecolor  = multGold,
  pdftitle   = {Phase Mirror: Adversarial Digital-Twin Framework},
  pdfauthor  = {Multiplicity Foundation},
  pdfsubject = {Defensive Prior-Art Publication},
}

%% ── Code Listings ───────────────────────────────────────────────────────────
\lstdefinestyle{pythonstyle}{
  backgroundcolor=\color{codeGray},
  commentstyle=\color{codeGreen}\itshape,
  keywordstyle=\color{codeBlue}\bfseries,
  stringstyle=\color{codeRed},
  basicstyle=\ttfamily\footnotesize,
  breaklines=true,
  captionpos=b,
  numbers=left,
  numberstyle=\tiny\color{gray},
  showspaces=false,
  tabsize=2,
  frame=single,
  rulecolor=\color{multBlue!40},
}
\lstdefinestyle{tsstyle}{
  backgroundcolor=\color{codeGray},
  commentstyle=\color{codeGreen}\itshape,
  keywordstyle=\color{codePurple}\bfseries,
  stringstyle=\color{codeRed},
  basicstyle=\ttfamily\footnotesize,
  breaklines=true,
  numbers=left,
  numberstyle=\tiny\color{gray},
  frame=single,
  rulecolor=\color{multBlue!40},
  morekeywords={interface,type,const,let,export,function,async,await,
                boolean,string,number,void,Promise},
}
\lstset{style=pythonstyle}

%% ── Theorem Environments ────────────────────────────────────────────────────
\theoremstyle{definition}
\newtheorem{definition}{Definition}[section]
\theoremstyle{plain}
\newtheorem{theorem}{Theorem}[section]
\newtheorem{lemma}[theorem]{Lemma}
\theoremstyle{remark}
\newtheorem{remark}{Remark}[section]
\newtheorem{invariant}{Invariant}[section]

%% ── Custom Boxes ────────────────────────────────────────────────────────────
\tcbuselibrary{skins,breakable}
\newtcolorbox{claimbox}[1][]{
  enhanced, title={Prior-Art Claim},
  colback=multBlue!5, colframe=multBlue,
  fonttitle=\bfseries, breakable, #1}
\newtcolorbox{alertbox}[1][]{
  enhanced, title={Governance Invariant (L0)},
  colback=codeRed!5, colframe=codeRed,
  fonttitle=\bfseries\color{codeRed}, breakable, #1}

%% ── Headers / Footers ───────────────────────────────────────────────────────
\pagestyle{fancy}
\fancyhf{}
\fancyhead[L]{\small\textcolor{multBlue}{\textbf{Phase Mirror} --- Adversarial Digital-Twin Framework}}
\fancyhead[R]{\small\textcolor{multGold}{Citizen Gardens}}
\fancyfoot[C]{\thepage}
\renewcommand{\headrulewidth}{0.4pt}
\setlength{\headheight}{15pt}

%% ── Section Formatting ──────────────────────────────────────────────────────
\titleformat{\section}{\large\bfseries\color{multBlue}}{{\thesection}}{1em}{}[\titlerule]
\titleformat{\subsection}{\normalsize\bfseries\color{multBlue!80}}{\thesubsection}{1em}{}
\titleformat{\subsubsection}{\normalsize\bfseries\color{multBlue!60}}{\thesubsubsection}{1em}{}

%% ── Document ────────────────────────────────────────────────────────────────
\begin{document}

%% ─────────────────────────────── TITLE PAGE ────────────────────────────────
\begin{titlepage}
\begin{center}
\vspace*{2cm}
{\Huge\bfseries\color{multBlue}Phase Mirror\\[0.4em]Adversarial Digital-Twin Framework}\\[1.2em]
{\large\color{multGold}\bfseries Comprehensive Technical Report \&\\Defensive Prior-Art Publication}\\[2.5em]
\hrule\vspace{0.6em}
{\normalsize\textbf{Citizen Gardens \\ The Foundation of Multiplicity}\\
\texttt{github.com/MultiplicityFoundation/PhaseMirror-HQ}}
\vspace{0.4em}\hrule\\[2em]
\begin{tabular}{ll}
\textbf{Document type:} & Defensive publication / prior-art disclosure \\
\textbf{Status:}        & Phase 5 operational; Phase 6 in progress \\
\textbf{ADR anchors:}   & ADR-001--015, ADR-MCRM-010--070, ADR-AGI-001--005 \\
\textbf{Core repo:}     & \texttt{MultiplicityFoundation/PhaseMirror-HQ} \\
\textbf{Date:}          & April 29, 2026 \\
\textbf{Version:}       & 1.0.0 \\
\end{tabular}
\vfill
\textcolor{gray}{\small
This document constitutes a defensive publication intended to establish prior art\\
for the Phase Mirror adversarial digital-twin framework and its constituent\\
mathematical structures. All concepts herein are hereby dedicated to the public\\
record to prevent third-party monopolisation of this design space.
}
\end{center}
\end{titlepage}

%% ─────────────────────────── TABLE OF CONTENTS ─────────────────────────────
\setcounter{tocdepth}{3}
\tableofcontents
\newpage

%% ─────────────────────── SECTION 1: EXECUTIVE SUMMARY ──────────────────────
\section{Executive Summary}
\label{sec:exec}

Phase Mirror is a mathematically governed, adversarial digital-twin framework
developed by the Citizen Gardens. It instantiates a \emph{prime-indexed,
recursively stable} substrate for modelling nonlinear, emergent structures in
computation, cognition, and physical systems. The framework is not a static
knowledge base; it is a living architecture in which every component---operators,
security layers, governance primitives, and runtime kernels---is recursively
derived from prime-number decompositions and verifiable contractivity invariants.

\subsection{Core Innovation Statement}

The central novelty of Phase Mirror is the fusion of three historically separate
disciplines into a single, self-consistent kernel:

\begin{enumerate}[label=(\roman*)]
  \item \textbf{Multiplicity Theory} --- a prime-indexed categorical foundation
    in which mathematical objects are reinterpreted as recursively generated
    interaction patterns governed by the core multiplicity operator
    $\Omega : \bigoplus_{p \in \mathbb{P}} \mathcal{H}_p \to \mathcal{H}$.
  \item \textbf{Adversarial Digital-Twin Simulation} --- a framework of
    dissonance detection, convergent causal resolution, and phase-weighted
    feedback loops that model and correct deviations in agentic systems.
  \item \textbf{Prime-Anchored Governance} --- an executable policy layer
    (ADRs, CI gates, Lean4 proofs, BLS12-381 attestation) that treats
    governance decisions as first-class mathematical objects.
\end{enumerate}

\subsection{Key Results to Date}

\begin{itemize}
  \item \textbf{Phase 5 (LLVM/MLIR)} operational: prime-indexed lowering pipeline
    with full ACE/PETC contractivity enforcement across 10{,}000-step simulations.
  \item \textbf{27+ ADRs} formally accepted, covering runtime kernel, QAGI
    governance, self-modification boundaries, BLS12-381 sovereign signing,
    and Microsoft AGT integration.
  \item \textbf{Lean4 proof pipeline} (ADR-014) scoped and under active
    formalisation: contractivity, ACE budget conservation, and Lobian guard
    properties are targeted for machine-checked proofs.
  \item \textbf{Prime-Weighted Execution Hashing (PWEH)} defined as the primary
    attestation layer binding execution steps to governance metadata in a
    post-quantum-safe hash chain.
  \item \textbf{QAGI stack} (ADR-010--013) formally governed: Prime-Indexed
    Tensor Networks, Non-Associative Fractal Learning, and Bayesian Quantum
    Network extensions each have dedicated acceptance criteria and gating.
\end{itemize}

\subsection{Purpose of This Document}

This report serves four purposes simultaneously:
\begin{enumerate}
  \item \textbf{Internal technical record} --- a comprehensive survey of current
    system architecture for engineering and research teams.
  \item \textbf{Defensive publication} --- a timestamped prior-art disclosure
    preventing third-party patents on any sub-component herein.
  \item \textbf{Investor / partner briefing} --- an accessible overview of the
    framework's unique value proposition and maturity level.
  \item \textbf{Academic precursor} --- a citable technical document preceding
    peer-reviewed publications on PITN, PWEH, NAFL, and related sub-systems.
\end{enumerate}

\newpage

%% ─────────────────── SECTION 2: SYSTEM ARCHITECTURE ────────────────────────
\section{System Architecture Overview}
\label{sec:arch}

\subsection{Layered Architecture}

Phase Mirror is structured as five co-dependent strata. The layers interact
through explicitly governed interfaces; no layer may reach across an
intermediate layer without an accepted ADR authorising that dependency.

\begin{table}[h!]
\centering
\caption{Phase Mirror Architectural Strata}
\begin{tabular}{clll}
\toprule
\textbf{Layer} & \textbf{Name} & \textbf{Key components} & \textbf{ADR anchor} \\
\midrule
L5 & QAGI Cognitive   & NAFL, BQN, PIRL, PEQZEA            & ADR-010--013 \\
L4 & Adversarial Twin & PMD, CCRE, PWCFL, surface levers    & ADR-002--004 \\
L3 & Governance       & AGT, CI gates, Lean4 proofs         & ADR-015, AGI-001--005 \\
L2 & Cryptographic    & BLS12-381, PWEH, HKDF-SHA256        & ADR-005, ADR-011 \\
L1 & Math Kernel      & $\Omega$, PIRTM, ACE, PETC          & ADR-001, ADR-009 \\
\bottomrule
\end{tabular}
\end{table}

\begin{mdframed}[linecolor=multBlue,linewidth=1pt]
\begin{center}
\textbf{Phase Mirror Layered Stack (schematic)}\\[0.4em]
\begin{tabular}{|c|}
\hline
L5 -- QAGI Cognitive Architecture (NAFL / BQN / PIRL / PEQZEA) \\
\hline
L4 -- Adversarial Digital-Twin (PMD / CCRE / PWCFL) \\
\hline
L3 -- Governance / CI (ADRs / Lean4 / AGT / SLO) \\
\hline
L2 -- Cryptographic Attestation (BLS12-381 / PWEH) \\
\hline
L1 -- Mathematical Kernel ($\Omega$ / PIRTM v2 / ACE / PETC) \\
\hline
\end{tabular}
\end{center}
\end{mdframed}

\subsection{L0 Invariants}

\begin{alertbox}
Three fields are absolutely protected across all layers: \texttt{prime\_index},
\texttt{epsilon}, and \texttt{op\_norm\_T}. No code path may modify these
without routing through the \texttt{lobian\_guard.py} contractivity pre-check
and an ADR-level amendment. Violations are hard aborts in CI.
\end{alertbox}

\subsection{Repository Structure (PhaseMirror-HQ)}

The \texttt{MultiplicityFoundation/PhaseMirror-HQ} repository contains:
\begin{itemize}
  \item \texttt{packages/multiplicity/core/} --- $\Omega$ operator, sovereign
    topos, Langlands registrar, lambda-M protocols, formal verification engine.
  \item \texttt{governance/} --- ADRs, self-modification boundary, kill switch,
    Lobian guard, rollback executor, schema migrations.
  \item \texttt{packages/lambda/apps/mcp-server/} --- PMD detection tools,
    CCRE resolver, PWCFL tracker, ACE certifier.
  \item \texttt{pirtm/} --- Prime-Indexed Recursive Tensor Machine v2 kernel.
  \item \texttt{test-harness/} --- architecture conformance, kill-switch,
    ACE conservation, and PETC boundary tests.
  \item \texttt{docs/adr/} --- 70+ architecture decision records.
\end{itemize}

\newpage

%% ─────────────────── SECTION 3: MATHEMATICAL FOUNDATIONS ───────────────────
\section{Mathematical Foundations}
\label{sec:math}

\subsection{Multiplicity Theory --- Core Principles}

Multiplicity Theory treats every mathematical object not as a static entity but
as a \emph{recursively generated interaction pattern} labelled by prime
decompositions.

\begin{definition}[Prime-Graded Hilbert Bundle]
\label{def:hilbert-bundle}
The prime-graded Hilbert bundle is the direct sum
\begin{equation}
  \mathcal{H} \;=\; \bigoplus_{p \in \mathbb{P}} \mathcal{H}_p,
\end{equation}
where $\mathbb{P}$ denotes the set of rational primes and each fibre
$\mathcal{H}_p$ carries a Banach-space norm $\|\cdot\|_p$. A state vector
$\psi \in \mathcal{H}$ is called a \emph{multiplicity state}; its support
$\operatorname{supp}(\psi) = \{p \in \mathbb{P} : \psi_p \neq 0\}$ is the
\emph{prime signature} of $\psi$.
\end{definition}

\begin{definition}[Core Multiplicity Operator $\Omega$]
\label{def:omega}
The core multiplicity operator is a specialisation functor
\begin{equation}
  \Omega : \mathcal{H} \;\longrightarrow\; \mathcal{H}, \qquad
  \Omega = \prod_{p \in \mathbb{P}} \Pi_p \circ \Phi_p,
\end{equation}
where $\Pi_p$ is the lawfulness projector onto $\mathcal{H}_p$ and $\Phi_p$ is
a sector-specific map (meta-relativity, Moonshine, CEQG, neuromorphic UME)
derived by calling \texttt{specialize(sector, **params)}.
Every existing module in \texttt{multiplicity.core} is a restriction of $\Omega$
to a single sector.
\end{definition}

\subsection{Multiplicity Category}

\begin{definition}[Multiplicity Category $\mathbf{Mult}$]
The Multiplicity Category has:
\begin{itemize}
  \item \textbf{Objects}: prime-graded Hilbert bundles $(\mathcal{H},\Pi)$
    equipped with ACE-budget functions $\alpha : \mathbb{P} \to \mathbb{R}_{>0}$.
  \item \textbf{Morphisms}: ACE-constrained, PETC-preserving maps
    $f : (\mathcal{H},\Pi,\alpha) \to (\mathcal{H}',\Pi',\alpha')$ such that
    $\|f\|_{\mathrm{op}} \le 1$ (contractivity) and
    $\mathrm{PETC}(f(\psi)) = \mathrm{PETC}(\psi)$ for all $\psi$.
\end{itemize}
Models in disparate domains (Moonshine, CEQG, Neuromorphic UME) are conjectured
to be faithful functors $F : \mathbf{Sect} \to \mathbf{Mult}$.
\end{definition}

\subsection{Prime-Indexed Recursive Tensor Machine (PIRTM v2)}

The PIRTM v2 is the computational realisation of the $\Omega$ operator at
runtime. Its fundamental update law is:

\begin{equation}
  \label{eq:pirtm-update}
  T_{t+1} \;=\; P\, F_t\, K_t(T_t),
\end{equation}
where $T_t \in \mathcal{H}$ is the state tensor at step $t$, $K_t$ is a
prime-indexed kernel, $F_t$ is an environment-indexed forgetfulness operator,
and $P$ is the ACE-budget projector.

\begin{invariant}[ACE Conservation]
\label{inv:ace}
For all steps $t$:
\begin{equation}
  \sum_{p \in \operatorname{supp}(T_{t+1})} \alpha(p) \;\le\;
  \sum_{p \in \operatorname{supp}(T_t)} \alpha(p).
\end{equation}
Any execution that violates this inequality is an L0 violation and must abort.
\end{invariant}

\begin{invariant}[PETC Environment Boundary]
\label{inv:petc}
Let $G$ and $H$ be two distinct prime environments. The composition
\begin{equation}
  \mathrm{Step}_{G,H}(T) \;\coloneqq\;
  \mathrm{Step}_H \circ \mathrm{Step}_G (T)
\end{equation}
is \textbf{forbidden}. Cross-environment PETC composition must be rejected at
both compile-time (lowering pipeline) and runtime (spectral gate).
\end{invariant}

\subsection{Phase-Weighted Contractivity Feedback Loop (PWCFL)}

The PWCFL monitors agent convergence on lever branches in real time. Let
$\delta_i = D_i - D_{i-1}$ be the dissonance delta at commit $i$. The
weighted average contraction (WAC) over a sliding window of size $N$ is:

\begin{equation}
  \label{eq:wac}
  \mathrm{WAC}_N \;=\;
  \frac{\displaystyle\sum_{i=t-N}^{t} e^{-\lambda(t-i)}\,|\delta_i|}
       {\displaystyle\sum_{i=t-N}^{t} e^{-\lambda(t-i)}},
  \qquad \lambda = 0.3 \text{ (default)}.
\end{equation}

\begin{invariant}[PWCFL Halt Condition]
If $\mathrm{WAC}_N > 1.0$ and $\mathrm{WAC}_N$ is monotonically increasing
over the window, the commit loop halts immediately and the lever state
transitions to \texttt{stalled}. No further commits may be made until a
human resolves the stall.
\end{invariant}

\subsection{Prime-Indexed Tensor Networks (PITNs)}

PITNs extend the PIRTM substrate to full tensor network semantics. A
prime-indexed tensor $\mathcal{T}$ is an element of the product space:

\begin{equation}
  \mathcal{T} \;\in\; \bigotimes_{p \in S} \mathcal{H}_p,
  \quad S \subset \mathbb{P} \text{ finite.}
\end{equation}

The \emph{non-associative} update rule for a PITN is:

\begin{equation}
  \label{eq:pitn-update}
  \mathcal{T}^{(t+1)} \;=\;
  \Phi_{p_\sigma(t)}\!\left(\mathcal{T}^{(t)}\right),
\end{equation}
where the result depends irreducibly on the \emph{path} (ordering) of prime
moves $\sigma$. This path-dependence is the foundation of Non-Associative
Fractal Learning (NAFL).

\begin{definition}[Prime Spectral Convergence]
A PITN trajectory $\{\mathcal{T}^{(t)}\}$ exhibits prime spectral convergence
if, for all $p \in S$, the $p$-sector projection satisfies:
\begin{equation}
  \lim_{t\to\infty} \|\Pi_p\,\mathcal{T}^{(t)} - \Pi_p\,\mathcal{T}^*\|_p
  \;=\; 0,
\end{equation}
where $\mathcal{T}^*$ is the fixed-point tensor determined by $\Omega$.
\end{definition}

\subsection{Prime-Weighted Execution Hashing (PWEH)}

PWEH binds execution steps, prime-indexed operator selections, tensor-state
observables, and governance metadata into a single integrity chain.

\begin{definition}[PWEH Integrity Chain]
Let $s_0, s_1, \ldots, s_n$ be the sequence of execution steps. The PWEH
chain is:
\begin{equation}
  \label{eq:pweh}
  H_0 \;\coloneqq\; \mathrm{H}(s_0 \,\|\, m_0), \qquad
  H_{k+1} \;\coloneqq\; \mathrm{H}\!\left(H_k \,\|\, s_{k+1} \,\|\, m_{k+1}
    \,\|\, \pi(k+1)\right),
\end{equation}
where $\pi(k)$ is the $k$-th rational prime, $m_k$ is the governance metadata
tuple at step $k$, and $\mathrm{H}$ is a post-quantum hash (HKDF-SHA256 or
ML-DSA-65 compatible).
\end{definition}

\begin{theorem}[PWEH Path Binding]
Any two execution histories $\sigma \neq \sigma'$ that differ in at least one
step-index / prime-label pair yield distinct PWEH chain termini
$H_n(\sigma) \neq H_n(\sigma')$ with overwhelming probability
(collision resistance of the underlying hash function).
\end{theorem}

\subsection{Non-Associative Fractal Learning (NAFL)}

NAFL is the cognitive architecture built above the PITN substrate. It
formalises the intuition that \emph{the history of prime moves constitutes the
effective state} and that knowledge is formed through ordered, irreversible
fractal refinements.

\begin{definition}[NAFL Decision Manifold]
Let $\mathcal{M}$ be a tensor manifold equipped with a prime-indexed
non-associative operator grammar $\mathcal{G} = \{g_p\}_{p \in \mathbb{P}}$.
The NAFL decision manifold is:
\begin{equation}
  \mathcal{D} \;\coloneqq\;
  \left\{ \gamma = (p_1, p_2, \ldots, p_k) \;\middle|\;
    p_i \in \mathbb{P},\;
    g_{p_k} \circ \cdots \circ g_{p_1}(\mathcal{T}_0) \in \mathcal{M}
  \right\}.
\end{equation}
Two paths $\gamma, \gamma'$ are \emph{equivalent} if and only if they terminate
at the same manifold point \textbf{and} have the same PWEH chain (no path
aliasing is permitted).
\end{definition}

\subsection{BLS12-381 Sovereign Signing}

All proof envelopes in Phase Mirror are authenticated by a BLS12-381 pairing
signature derived via a two-layer HKDF:

\begin{equation}
  k_{\text{circuit}} \;=\;
  \mathrm{HKDF\text{-}SHA256}\!\left(
    \text{salt} = \texttt{mtpi:sovereign-sign:v1},\;
    \text{ikm} = \text{identitySecret},\;
    \text{info} = \text{circuit}
  \right).
\end{equation}

\begin{alertbox}[title={L0 Key Boundary Invariant}]
No circuit runtime in a server-side, relay, or cloud context shall ingest raw
\texttt{identitySecret}. Only HKDF-derived keys may cross this boundary.
Violations return \texttt{HTTP 400 / ERR\_SECRET\_BOUNDARY}.
\end{alertbox}

\newpage

%% ─────────────────────── SECTION 4: ADVERSARIAL TWIN ───────────────────────
\section{Adversarial Digital-Twin Framework}
\label{sec:adt}

\subsection{Phase Mirror Dissonance (PMD)}

The adversarial digital-twin operates through the \emph{Phase Mirror Dissonance}
engine, which continuously scans the repository for deviations between the
mathematical ideal (ADR governance contracts, L0 invariants) and the actual
implementation state.

PMD findings are classified by governance tier:

\begin{table}[h!]
\centering
\caption{PMD Governance Tiers}
\begin{tabular}{cllc}
\toprule
\textbf{Tier} & \textbf{Finding type} & \textbf{Typical cause} & \textbf{Tension} \\
\midrule
\texttt{pass}  & Clean state          & No deviation detected           & 0.0 \\
\texttt{warn}  & \texttt{drift}       & Implementation diverged from ADR & 0.5 \\
\texttt{block} & \texttt{l0-violation}& L0-protected field exposed       & 1.0 \\
\texttt{block} & \texttt{adr-gap}     & Decision made without ADR anchor & 1.0 \\
\texttt{block} & \texttt{consent}     & Data-path change without consent & 1.0 \\
\bottomrule
\end{tabular}
\end{table}

\subsection{Convergent Causal Resolution Engine (CCRE)}

The CCRE is the ordering layer between detection (PMD) and action
(lever branches). It applies three ordering rules in strict priority:

\begin{enumerate}
  \item \textbf{ADR-before-code}: any \texttt{adr-gap} lever precedes any
    \texttt{drift} lever on the same component.
  \item \textbf{Consent-before-data}: any \texttt{consent} lever precedes
    any \texttt{drift} lever on data-handling paths (PMD L0-003).
  \item \textbf{Tension-descending}: among unordered levers,
    highest tension is resolved first.
\end{enumerate}

The CCRE builds a dependency DAG and topologically sorts it. A cycle in the
DAG is a \textbf{governance conflict} requiring human intervention.

\begin{equation}
  \mathrm{DAG} = (V, E), \quad
  (A \to B) \in E \;\Leftrightarrow\;
  \text{ordering rule places } A \text{ before } B.
\end{equation}

\subsection{Three MCP Governance Tools (ADR-002)}

\begin{table}[h!]
\centering
\caption{Phase Mirror MCP Tool Loop}
\begin{tabular}{llp{6.5cm}}
\toprule
\textbf{Tool} & \textbf{Phase} & \textbf{Action} \\
\midrule
\texttt{surface\_levers}          & Detect  & Run PMD; group findings into ranked levers with CCRE ordering. \\
\texttt{create\_lever\_branch}    & Act     & Create \texttt{lever/*} branch; enforce CCRE dependency; start PWCFL monitoring. \\
\texttt{certify\_lever\_convergence} & Certify & Run ACE WAC telescoping on full branch history; return binary verdict. \\
\bottomrule
\end{tabular}
\end{table}

\subsection{Self-Modification Boundary (ADR-AGI-004)}

Seven files in \texttt{governance/self\_modification/} are subject to binding
rules. The key constraint for each:

\begin{itemize}
  \item \texttt{proposal.py} --- may not propose changes to L0-protected fields
    without an ADR-level amendment on a live session.
  \item \texttt{daemon.py} --- must call \texttt{validation\_gates.py} before
    any proposal execution; must call \texttt{lobian\_guard.py} before any
    change touching spectral parameters.
  \item \texttt{watchdog.py} --- must emit a ledger entry for every state
    transition; silent state changes are invariant violations.
  \item \texttt{kill\_switch.py} --- trigger condition must have a passing CI
    test; a kill switch with no trigger test is treated as absent.
  \item \texttt{lobian\_guard.py} --- must perform contractivity pre-check
    before approving any spectral parameter change.
  \item \texttt{rollback\_executor.py} --- must verify ledger quorum before
    executing any rollback.
  \item \texttt{schema\_migrations.py} --- must not alter L0-protected column
    semantics without an ADR amendment.
\end{itemize}

\newpage

%% ────────────────────────── SECTION 5: CODE SNIPPETS ───────────────────────
\section{Representative Code Artefacts}
\label{sec:code}

\subsection{PIRTM v2 Step with ACE Budget Enforcement}

\begin{lstlisting}[language=Python, style=pythonstyle,
  caption={PIRTM v2 core update step --- Eq.~(\ref{eq:pirtm-update})},
  label={lst:pirtm}]
# pirtm/core/pirtm_v2.py
import numpy as np
from typing import Callable

class L0InvariantViolation(RuntimeError):
    """Raised when an L0-protected invariant is violated."""

class PIRTMv2:
    """
    Prime-Indexed Recursive Tensor Machine, version 2.
    Implements T_{t+1} = P F_t K_t(T_t) with ACE/PETC enforcement.
    """
    def __init__(self, prime_index: dict,
                 epsilon: float, op_norm_T: float):
        # L0-protected fields -- only writable via lobian_guard
        self._prime_index = prime_index
        self._epsilon = epsilon
        self._op_norm_T = op_norm_T

    def step(self, T: np.ndarray, kernel: Callable,
             forget: Callable, ace_budget: dict, env_id: str
             ) -> np.ndarray:
        """Single PIRTM step with L0-invariant enforcement."""
        T_k   = kernel(T, self._prime_index)     # 1. Apply K_t
        T_f   = forget(T_k, env_id)              # 2. Apply F_t
        T_nxt = self._ace_project(T_f, ace_budget)  # 3. Project P

        # 4. Contractivity check (L0)
        gain = (self._spectral_radius(T_nxt)
                / (self._spectral_radius(T) + 1e-12))
        if gain >= self._op_norm_T + self._epsilon:
            raise L0InvariantViolation(
                f"Contractivity violated: gain={gain:.4f}")
        return T_nxt

    def _ace_project(self, T, budget):
        before = sum(budget.get(p, 0.0)
                     for p in self._prime_index)
        proj = T.copy()
        after = sum(budget.get(p, 0.0)
                    * np.linalg.norm(proj[..., i])
                    for i, p in enumerate(self._prime_index))
        if after > before:
            proj *= before / (after + 1e-12)
        return proj

    def _spectral_radius(self, T):
        return float(np.linalg.norm(T, ord=2))
\end{lstlisting}

\subsection{CCRE Resolution Queue Builder (TypeScript)}

\begin{lstlisting}[style=tsstyle,
  caption={CCRE dependency DAG construction and topological sort},
  label={lst:ccre}]
// src/utils/ccre-resolver.ts
export function buildResolutionQueue(
  findings: PMDFinding[]
): ResolutionQueue | CcreConflict {

  const entries = findings.map((f, i) => ({
    lever_id: `${f.type}-${String(i).padStart(2,"0")}`,
    title: f.title,
    priority: 0,
    blocks: [] as string[],
    state: "pending" as const,
    rationale: "",
  }));

  const edges: [string, string][] = [];
  for (let a = 0; a < entries.length; a++) {
    for (let b = 0; b < entries.length; b++) {
      if (a === b) continue;
      const fa = findings[a], fb = findings[b];

      // Rule 1: ADR-before-code on same component
      if (fa.type === "adr-gap" && fb.type === "drift"
          && fa.component === fb.component) {
        edges.push([entries[a].lever_id, entries[b].lever_id]);
        entries[a].rationale =
          "Architecture decision must precede implementation.";
      }
      // Rule 2: Consent-before-data (L0-003)
      if (fa.type === "consent" && fb.type === "drift"
          && fb.path_type === "data") {
        edges.push([entries[a].lever_id, entries[b].lever_id]);
        entries[a].rationale =
          "Consent must precede data-processing changes.";
      }
    }
  }

  for (const [from, to] of edges) {
    const e = entries.find(e => e.lever_id === from)!;
    if (!e.blocks.includes(to)) e.blocks.push(to);
  }

  const cycle = detectCycle(entries, edges);
  if (cycle) return { type: "CcreConflict", cycle };

  const sorted = topologicalSort(entries, edges);
  sorted.sort((a, b) => {
    const tA = findings[entries.indexOf(a)].tension ?? 0;
    const tB = findings[entries.indexOf(b)].tension ?? 0;
    return tB - tA;           // tension-descending
  });
  sorted.forEach((e, i) => { e.priority = i + 1; });
  return { ordered: sorted };
}
\end{lstlisting}

\subsection{PWEH Integrity Chain Construction}

\begin{lstlisting}[language=Python, style=pythonstyle,
  caption={PWEH chain construction --- Eq.~(\ref{eq:pweh})},
  label={lst:pweh}]
# pirtm/security/pweh.py
import hashlib, hmac, struct, json
from sympy import nextprime

_PRIMES: list[int] = []

def _prime(k: int) -> int:
    """Return the k-th rational prime (1-indexed)."""
    while len(_PRIMES) < k:
        base = _PRIMES[-1] if _PRIMES else 1
        _PRIMES.append(nextprime(base))
    return _PRIMES[k - 1]

def build_pweh_chain(steps: list[bytes],
                     metadata: list[dict],
                     salt: bytes = b"pweh:v1"
                     ) -> list[bytes]:
    """
    H_0 = H(step_0 || meta_0)
    H_{k+1} = H(H_k || step_{k+1} || meta_{k+1} || prime(k+1))
    """
    assert len(steps) == len(metadata)
    chain: list[bytes] = []
    for k, (s, m) in enumerate(zip(steps, metadata)):
        prime_b = struct.pack(">Q", _prime(k + 1))
        meta_b  = json.dumps(m, sort_keys=True,
                             separators=(",",":")).encode()
        payload = (chain[-1] + s + meta_b + prime_b
                   if k > 0 else s + meta_b)
        chain.append(
            hmac.new(salt, payload, hashlib.sha256).digest())
    return chain

def verify_pweh_chain(chain, steps, metadata,
                      salt=b"pweh:v1") -> bool:
    expected = build_pweh_chain(steps, metadata, salt)
    return all(hmac.compare_digest(a, b)
               for a, b in zip(chain, expected))
\end{lstlisting}

\subsection{NAFL Non-Associative Decision Manifold}

\begin{lstlisting}[language=Python, style=pythonstyle,
  caption={NAFL path-dependent prime-ordered operator application},
  label={lst:nafl}]
# qagi/nafl/decision_manifold.py
from dataclasses import dataclass, field
from typing import Sequence
import numpy as np

@dataclass
class NAFLState:
    tensor:      np.ndarray
    path:        list[int]   = field(default_factory=list)
    pweh_chain:  list[bytes] = field(default_factory=list)

class NAFLDecisionManifold:
    """
    g_{p_k} o ... o g_{p_1}(T_0) depends on the ordering of primes.
    Different orderings yield different knowledge-manifold points.
    """
    def __init__(self, operators: dict[int, callable]):
        self.operators = operators      # {prime -> operator}

    def apply_path(self, T0: np.ndarray,
                   path: Sequence[int]) -> NAFLState:
        state = NAFLState(tensor=T0.copy())
        for p in path:
            if p not in self.operators:
                raise ValueError(f"No operator for prime {p}")
            state.tensor = self.operators[p](state.tensor)
            state.path.append(p)
            state.pweh_chain.append(state.tensor.tobytes()[:32])
        return state

    def paths_equivalent(self, s1: NAFLState,
                          s2: NAFLState, tol=1e-8) -> bool:
        """Equivalent IFF same terminal point AND same PWEH chain."""
        return (np.allclose(s1.tensor, s2.tensor, atol=tol)
                and s1.pweh_chain == s2.pweh_chain)
\end{lstlisting}

\subsection{Lobian Guard Contractivity Pre-Check}

\begin{lstlisting}[language=Python, style=pythonstyle,
  caption={Lobian guard protecting L0 spectral fields},
  label={lst:lobian}]
# governance/self_modification/lobian_guard.py
from dataclasses import dataclass

@dataclass
class SpectralParams:
    prime_index: dict
    epsilon:     float
    op_norm_T:   float

class LobianGuard:
    """
    Guards ALL modifications to L0-protected spectral parameters.
    A failed pre-check is a hard abort.
    """
    def pre_check(self, current: SpectralParams,
                  proposed: SpectralParams) -> None:
        if proposed.op_norm_T >= 1.0:
            raise L0InvariantViolation(
                f"op_norm_T={proposed.op_norm_T} >= 1.0 "
                "violates contractivity")
        if proposed.epsilon < 0:
            raise L0InvariantViolation(
                f"epsilon={proposed.epsilon} < 0 is non-physical")
        missing = (set(current.prime_index)
                   - set(proposed.prime_index))
        if missing:
            raise L0InvariantViolation(
                f"Silent prime_index removal: {missing}")

class L0InvariantViolation(RuntimeError):
    """Raised on any L0-protected invariant violation."""
\end{lstlisting}

\newpage

%% ──────────────────── SECTION 6: GOVERNANCE ARCHITECTURE ───────────────────
\section{Governance Architecture}
\label{sec:gov}

\subsection{ADR Hierarchy}

\begin{table}[h!]
\centering
\caption{ADR Namespace Structure (April 2026)}
\begin{tabular}{lll}
\toprule
\textbf{Namespace} & \textbf{Scope} & \textbf{Count} \\
\midrule
ADR-001--015     & Core kernel, tools, QAGI, AGT   & 15 \\
ADR-MCRM-010--070 & Multiplicity Core Roadmap       & 7 \\
ADR-AGI-001--005 & AGI safety \& contractivity     & 5 \\
ADR-MVP-001--003 & MVP schema \& transport         & 3 \\
ADR-Phase5--7    & LLVM/MLIR phases                & 6 \\
ADR-PIRDS-001--002 & Prime-Indexed Runtime DS      & 2 \\
\bottomrule
\end{tabular}
\end{table}

\subsection{CI Gate Sequencing}

Every PR merging into \texttt{main} must pass five CI gates in strict order:

\begin{enumerate}
  \item \textbf{ADR compliance gate} (\texttt{adr-gate.yml}) --- verifies that
    every changed module has a corresponding accepted ADR.
  \item \textbf{MCP surface scan} (\texttt{mcp-scan.yml}) --- deterministic
    security scanning for all \texttt{mcp/} endpoints.
  \item \textbf{Policy coverage} (\texttt{policy-coverage.yml}) --- verifies
    machine-readable YAML policies are in sync with constitutional documents.
  \item \textbf{AGT verification} (\texttt{agt-verify.yml}) --- Microsoft AGT
    content validation and SRE-grade observability checks.
  \item \textbf{Kill-switch test} (\texttt{test\_r04\_kill\_switch.py}) ---
    armed and disarmed states both pass before any merge.
\end{enumerate}

\subsection{Microsoft AGT Integration (ADR-015)}

Phase Mirror adopted the Microsoft Agentic Graph Transformer (AGT) integration
roadmap across four phases, all now marked \textbf{COMPLETE}:

\begin{table}[h!]
\centering
\caption{AGT Integration Phase Status}
\begin{tabular}{clll}
\toprule
\textbf{Phase} & \textbf{Name} & \textbf{Key artefacts} & \textbf{Status} \\
\midrule
0 & Foundation     & ADR infrastructure, CI gates         & Complete \\
1 & MCP Hardening  & \texttt{adr-009}, \texttt{adr-010}   & Complete \\
2 & Policy-as-Code & \texttt{adr-011}, \texttt{governance/policies/} & Complete \\
3 & SRE + Identity & \texttt{adr-012}, \texttt{adr-013}, ML-DSA-65  & Complete \\
\bottomrule
\end{tabular}
\end{table}

\newpage

%% ─────────────────────────── SECTION 7: QAGI STACK ─────────────────────────
\section{Quantum-AGI Stack (ADR-010--013)}
\label{sec:qagi}

\subsection{Stack Layers}

\begin{table}[h!]
\centering
\caption{QAGI Sub-System Governance}
\begin{tabular}{lllc}
\toprule
\textbf{Layer} & \textbf{Sub-system} & \textbf{Key concept} & \textbf{ADR} \\
\midrule
Substrate   & PITN                 & Prime-indexed non-associative tensor update   & 010 \\
Security    & PWEH                 & Post-quantum integrity chain                  & 011 \\
Cognition   & NAFL                 & Path-dependent hierarchical knowledge         & 012 \\
Extensions  & BQN / PIRL / PEQZEA  & Bayesian nets, swarm, Zeno stability          & 013 \\
\bottomrule
\end{tabular}
\end{table}

\subsection{Bayesian Quantum Network (BQN)}

BQN introduces probabilistic exploration into the NAFL decision manifold. A
\emph{tunneling policy} $\tau : \mathcal{D} \to [0,1]$ assigns a probability
of transitioning to a non-greedy prime path. The governance constraint is:
\begin{equation}
  \tau(\gamma) \;\le\; \tau_{\max}(|\gamma|), \qquad
  \tau_{\max}(k) = \frac{1}{1 + e^{\lambda(k - k_0)}},
\end{equation}
where $k = |\gamma|$ is the path length and $\lambda, k_0$ are policy
hyperparameters requiring signed ADR amendments to change.

\subsection{Prime-Swarm Intelligence Research Layer (PIRL)}

PIRL assigns each prime $p \in S$ a \emph{role agent} responsible for
exploring the $\mathcal{H}_p$ fibre. The prime role map
$\rho : \mathbb{P} \to \mathcal{R}$ (roles: \textsc{Explorer},
\textsc{Consolidator}, \textsc{Auditor}) is fixed at initialisation and may
only be mutated via a signed governance transition.

\subsection{Prime-Embedded Quantum Zeno Effect (PEQZEA)}

PEQZEA exploits the quantum Zeno effect to lock the PITN state into a
certified attractor. Let $P_\star$ be the projector onto a target attractor.
PEQZEA repeatedly applies:
\begin{equation}
  \mathcal{T}^{(t+1)} = P_\star\,\Phi_{p(t)}(\mathcal{T}^{(t)}),
\end{equation}
at a rate $\nu \gg 1/\tau_{\mathrm{decoherence}}$. An \emph{unlock policy}
(ADR-013 Issue I3) specifies when $P_\star$ may be relaxed.

\newpage

%% ──────────────── SECTION 8: DEFENSIVE PRIOR-ART CLAIMS ────────────────────
\section{Defensive Prior-Art Publication}
\label{sec:priorart}

This section constitutes the formal defensive publication. Each numbered claim
below identifies a specific technical concept first disclosed herein or in the
referenced GitHub repository at
\url{https://github.com/MultiplicityFoundation/PhaseMirror-HQ}.
Publication of this document establishes a prior-art date of
\textbf{April 29, 2026}.

\begin{claimbox}[title={Claim 1 --- Core Multiplicity Operator $\Omega$}]
A prime-graded Hilbert-bundle operator $\Omega$ that unifies disparate physical
and cognitive models (Moonshine, CEQG, Meta-Relativity, neuromorphic UME) as
sector-specific restrictions via a \texttt{specialize(sector, **params)} API,
where the unified framework preserves a single ACE-budget conservation invariant
across all sectors.
\end{claimbox}

\begin{claimbox}[title={Claim 2 --- PIRTM v2 Update Law with PETC Boundary}]
A recursive tensor machine update law $T_{t+1} = P F_t K_t(T_t)$ in which:
(a) $P$ is a prime-indexed ACE-budget projector enforcing energy conservation;
(b) $F_t$ is an environment-indexed forgetfulness operator;
(c) cross-environment composition $\mathrm{Step}_{G,H}$ is rejected at both
compile-time and runtime via PETC boundary checks.
\end{claimbox}

\begin{claimbox}[title={Claim 3 --- Phase Mirror Dissonance (PMD) Framework}]
An adversarial digital-twin governance system that: (a) scans repositories for
deviations between ADR-encoded mathematical ideals and implementation state;
(b) classifies findings as \texttt{warn} (tension 0.5) or \texttt{block}
(tension 1.0); (c) generates ranked lever manifests with CCRE dependency
ordering; (d) certifies convergence via WAC telescoping against an ACE budget.
\end{claimbox}

\begin{claimbox}[title={Claim 4 --- Convergent Causal Resolution Engine (CCRE)}]
An ordering engine for governance findings that builds a directed acyclic graph
from three precedence rules (ADR-before-code, consent-before-data,
tension-descending) and topologically sorts findings into a resolution queue,
with cycle detection returning a governance conflict rather than an autonomous
break.
\end{claimbox}

\begin{claimbox}[title={Claim 5 --- Phase-Weighted Contractivity Feedback Loop (PWCFL)}]
A per-commit convergence monitor embedded in the agent commit loop that
computes an exponentially-weighted average contraction (WAC) over a sliding
window and halts agent work when $\mathrm{WAC}_N > 1.0$ is monotonically
increasing, preventing runaway commit loops in autonomous agentic systems.
\end{claimbox}
\newpage
\begin{claimbox}[title={Claim 6 --- Prime-Weighted Execution Hashing (PWEH)}]
A post-quantum integrity chain in which each execution step's hash incorporates
the $k$-th rational prime as a path label:
$H_{k+1} = \mathrm{H}(H_k \| s_{k+1} \| m_{k+1} \| \pi(k+1))$,
binding execution history, governance metadata, and prime-indexed operator
selections into a single unforgeable attestation.
\end{claimbox}

\begin{claimbox}[title={Claim 7 --- Non-Associative Fractal Learning (NAFL)}]
A cognitive architecture over a PITN substrate in which: (a) the ordered
history of prime-indexed operator applications constitutes the effective
knowledge state; (b) different orderings of the same prime set yield distinct
manifold points; (c) two paths are equivalent only if they terminate at the
same point \textbf{and} have identical PWEH chains, preventing path aliasing.
\end{claimbox}

\begin{claimbox}[title={Claim 8 --- BLS12-381 Sovereign Signing with HKDF Isolation}]
A proof-envelope signing model in which: (a) the raw \texttt{identitySecret}
never crosses a relay boundary; (b) a circuit-specific signing key is derived
via HKDF-SHA256 with salt \texttt{mtpi:sovereign-sign:v1}; (c) relay-level
submissions containing raw identity secrets are rejected with
\texttt{ERR\_SECRET\_BOUNDARY}.
\end{claimbox}

\begin{claimbox}[title={Claim 9 --- Lobian Guard with L0 Spectral Pre-Check}]
A governance component (\texttt{lobian\_guard.py}) that enforces: (a) no
spectral parameter change (\texttt{prime\_index}, \texttt{epsilon},
\texttt{op\_norm\_T}) may proceed if $\texttt{op\_norm\_T} \ge 1.0$;
(b) silent removal of prime-index entries is a hard abort; (c) all seven files
in the self-modification surface are required to route through this guard
before touching L0-protected fields.
\end{claimbox}

\begin{claimbox}[title={Claim 10 --- PEQZEA Attractor Locking}]
A stabilisation technique for PITN trajectories that applies repeated attractor
projection $P_\star$ at a rate exceeding the decoherence timescale, governed by
a signed unlock policy requiring council approval before the attractor lock is
relaxed.
\end{claimbox}

\newpage

%% ──────────────────── SECTION 9: DEVELOPMENT ROADMAP ───────────────────────
\section{Development Roadmap}
\label{sec:roadmap}

\begin{table}[h!]
\centering
\caption{Phased Kernel + QAGI Roadmap (ADR-009)}
\begin{tabular}{p{1.3cm}p{1.9cm}p{4.2cm}p{4.6cm}}
\toprule
\textbf{Phase} & \textbf{Horizon} & \textbf{Objective} & \textbf{Acceptance criteria} \\
\midrule
0 & Weeks 1--4   & Repo hygiene, installable package, Rank 2 kernel disclosure & \texttt{pip install -e .} passes; \texttt{v0.1-kernel} tagged with SHA-256 manifest \\
1 & Weeks 5--16  & Rank 1 \& 2 math + PIRTM runtime core & ACE+PETC 10k-step simulation with zero conservation failures \\
2 & Weeks 17--28 & Rank 3: CMT, ETP, ZSD, quantum algorithms & CMT conservation verified; ZSD entropy improvement documented \\
3 & Weeks 29--40 & Rank 4: socio-atomic modelling, reciprocity & Stable vs.\ exploitative regimes reproduced \\
4 & Weeks 41--52 & Rank 5: inversion, fractional multiplicity, anti-resonance & Formal theorem or counterexample published; anti-resonance monitor calibrated \\
\bottomrule
\end{tabular}
\end{table}

\newpage

%% ──────────────────── SECTION 10: TENSIONS \& PRECISION QS ─────────────────
\section{Active Tensions and Precision Questions}
\label{sec:tensions}

\begin{mdframed}[linecolor=multGold, linewidth=1.5pt]
\textbf{Tension 1 --- Accuracy vs.\ Compliance}\\[0.2em]
Are we optimising for mathematical completeness (full Lean4 coverage, all
sectors unified under $\Omega$) or for operational compliance (minimal proofs
that satisfy CI gates)?\\[0.3em]
\textbf{Lever}: Owner = Multiplicity Foundation;
Metric = Lean4 coverage \%;
Horizon = 90 days.
\end{mdframed}

\vspace{0.8em}
\begin{mdframed}[linecolor=multGold, linewidth=1.5pt]
\textbf{Tension 2 --- Autonomy vs.\ Governance}\\[0.2em]
Should NAFL/PIRL prime-role assignments be mutable under signed governance
transitions, or fixed-at-initialisation to prevent adversarial role drift?\\[0.3em]
\textbf{Lever}: Owner = AGI safety team;
Metric = Unauthorised role-map mutations in 30 days;
Horizon = 30 days.
\end{mdframed}

\vspace{0.8em}
\begin{mdframed}[linecolor=multGold, linewidth=1.5pt]
\textbf{Tension 3 --- Open-Core vs.\ Proprietary}\\[0.2em]
Which sub-systems (PWEH circuit parameters, PEQZEA unlock policies, BQN
tunneling bounds) should be published as open prior art vs.\ retained as
competitive moat?\\[0.3em]
\textbf{Lever}: Owner = Foundation executive;
Metric = Governed-to-disclosed artefact ratio;
Horizon = 7 days.
\end{mdframed}

\vspace{0.8em}
\begin{mdframed}[linecolor=multGold, linewidth=1.5pt]
\textbf{Tension 4 --- Velocity vs.\ Safety}\\[0.2em]
PWCFL's conservative WAC threshold ($\lambda = 0.3$) produces false-positive
stalls. Should the threshold be relaxed for faster iteration, risking divergent
commit loops?\\[0.3em]
\textbf{Lever}: Owner = DevOps team;
Metric = False-positive stall rate per 100 commits;
Horizon = 30 days.
\end{mdframed}

\newpage

%% ──────────────────────── SECTION 11: FUTURE WORK ──────────────────────────
\section{Recommended Future Work}
\label{sec:future}

\begin{enumerate}
  \item \textbf{Lean4 Contractivity Proof Sprint (90 days)}: Formalise the
    PIRTM v2 update law and Lobian guard in Lean4. This is the single
    highest-leverage proof investment; it directly enables machine-checked
    CI gates.

  \item \textbf{PWEH Post-Quantum Upgrade (30 days)}: Replace HMAC-SHA256
    with ML-DSA-65 (CRYSTALS-Dilithium) as the attestation primitive, in
    line with ADR-015 Phase 3 federation identity requirements.

  \item \textbf{NAFL Benchmark Suite (60 days)}: Define a concrete decision-
    trace benchmark dataset for NAFL path-dependent equivalence testing.
    Required to close ADR-012 Issue I1.

  \item \textbf{Socio-Atomic Simulation Release (90 days)}: Complete the Rank 4
    PIRTM phase (ADR-009 Phase 3), producing a reproducible socio-atomic
    dataset and reciprocity analytics notebook.

  \item \textbf{Anti-Resonance Monitor Calibration (60 days)}: Implement and
    calibrate the anti-resonance monitor against Beatty, Sturmian, Fibonacci,
    Thue-Morse, and synthetic-comb schedule families.

  \item \textbf{BQN Tunneling Bounds ADR (30 days)}: Close ADR-013 Issue I1
    by defining explicit governance bounds for BQN tunneling at the
    decision-type level.

  \item \textbf{Fractal Spectrum Mapping Schema (30 days)}: Define the Fractal
    Spectrum trace validity schema required by ADR-013 Issue I4, enabling
    explainability audits for all governed QAGI decision contexts.

  \item \textbf{PIRL Council Governance (60 days)}: Draft the process for
    PEQZEA unlock approvals and PIRL role-map transitions, closing ADR-013
    Issues I2 and I3.
\end{enumerate}

\newpage

%% ─────────────────────────── SECTION 12: CONCLUSION ────────────────────────
\section{Conclusion}
\label{sec:conclusions}

Phase Mirror represents a fundamentally novel approach to agentic AI governance:
rather than treating governance as a policy layer bolted on top of a working
system, it encodes governance \emph{within} the mathematical substrate.
Prime-indexed operators, ACE-budget conservation, PETC environment boundaries,
and PWEH attestation are not policy annotations---they are theorems the system
must satisfy to execute at all.

The adversarial digital-twin loop (PMD $\to$ CCRE $\to$ PWCFL $\to$ ACE)
closes the gap between specification and implementation by making deviations
automatically detectable, causally ordered, convergence-monitored, and
cryptographically certified. The QAGI stack (PITNs, PWEH, NAFL, BQN, PIRL,
PEQZEA) extends this substrate towards artificial general intelligence while
maintaining explicit governance boundaries at every capability level.

This document establishes a public prior-art record as of \textbf{April 29, 2026}
for all ten defensive claims in Section~\ref{sec:priorart}.

\vspace{1.5em}
\hrule
\vspace{0.8em}
\noindent\textbf{Repository:}
\url{https://github.com/MultiplicityFoundation/PhaseMirror-HQ}\\
\textbf{Governing ADRs:} ADR-001--015, ADR-MCRM-010--070, ADR-AGI-001--005\\
\textbf{Prior-art date:} April 29, 2026\\
\textbf{Publisher:} Multiplicity Foundation

\appendix

This appendix records the explicit mathematical statements, operator-norm bounds,
and proof sketches implied by the architectural decisions underlying the Phase
Mirror framework. Its purpose is not to claim final formal completeness, but to
make precise the contractive estimates, conservation statements, and path-binding
properties that the main article relies upon.

The appendix focuses on five objects:
\begin{enumerate}[label=(\roman*)]
  \item the core multiplicity operator $\Omega$ on the prime-graded Hilbert bundle,
  \item the PIRTM update operator $T \mapsto P F_t K_t(T)$,
  \item the ACE budget projector $P$,
  \item the PWCFL weighted contraction statistic,
  \item the PWEH path-binding chain.
\end{enumerate}

Throughout, we work with a prime-graded Hilbert bundle
\begin{equation}
\mathcal{H} = \bigoplus_{p\in\mathbb{P}} \mathcal{H}_p,
\end{equation}
endowed with the norm
\begin{equation}
\|x\|_{\mathcal{H}}^2 = \sum_{p\in\mathbb{P}} w_p \|x_p\|_p^2,
\qquad w_p > 0,
\end{equation}
with only finitely many nonzero coordinates for each state under runtime use.

\section{Standing Assumptions}

\begin{definition}[Sectorwise boundedness]
For each prime $p$, let $\Phi_p : \mathcal{H}_p \to \mathcal{H}_p$ be a bounded
operator with operator norm satisfying
\begin{equation}
\|\Phi_p\|_{\mathrm{op}} \le \kappa_p.
\end{equation}
We say the family is uniformly contractive if
\begin{equation}
\sup_{p\in\mathbb{P}} \kappa_p \le \kappa < 1.
\end{equation}
\end{definition}

\begin{definition}[Lawfulness projector]
For each prime sector $p$, let $\Pi_p : \mathcal{H}_p \to \mathcal{H}_p$ denote
an orthogonal projector, hence idempotent and self-adjoint, satisfying
\begin{equation}
\Pi_p^2 = \Pi_p, \qquad \Pi_p^* = \Pi_p.
\end{equation}
Consequently,
\begin{equation}
\|\Pi_p\|_{\mathrm{op}} = 1
\end{equation}
unless $\Pi_p = 0$.
\end{definition}

\begin{definition}[Environment-indexed forget operator]
At runtime step $t$, let $F_t : \mathcal{H} \to \mathcal{H}$ be bounded with
\begin{equation}
\|F_t\|_{\mathrm{op}} \le \mu_t.
\end{equation}
The contractive regime is the case $\mu_t \le 1$ for all $t$, while strict
forgetful contraction corresponds to $\mu_t \le \mu < 1$.
\end{definition}

\begin{definition}[ACE projector]
Let $P : \mathcal{H} \to \mathcal{H}$ be the ACE-budget projection acting by
radial rescaling or orthogonal projection onto the closed feasible set
\begin{equation}
\mathcal{A} = \left\{x\in\mathcal{H} : \sum_{p\in\operatorname{supp}(x)}
\alpha(p)\,\vartheta_p(x_p) \le B\right\},
\end{equation}
where $\alpha(p) > 0$ are prime weights and $\vartheta_p(\cdot)$ is a positive
sector observable, typically a norm or seminorm. In all cases used here,
$P$ is assumed nonexpansive:
\begin{equation}
\|P(x)-P(y)\|_{\mathcal{H}} \le \|x-y\|_{\mathcal{H}}.
\end{equation}
\end{definition}

\begin{invariantbox}
The nonexpansiveness of $P$ is the mathematical abstraction of ACE budget
conservation. In the concrete implementation by radial scaling, it is immediate
because scaling by a factor in $[0,1]$ cannot increase distance from the origin.
For orthogonal projection onto a closed convex feasible set, nonexpansiveness is
standard in Hilbert space geometry.
\end{invariantbox}

\section{Norm Bounds for the Core Operator $\Omega$}

\begin{definition}[Core operator]
Define the sectorwise core multiplicity operator by
\begin{equation}
\Omega x = \bigoplus_{p\in\mathbb{P}} \Pi_p \Phi_p x_p.
\end{equation}
\end{definition}

\begin{lemma}[Sectorwise operator bound]
For every prime $p$,
\begin{equation}
\|\Pi_p \Phi_p\|_{\mathrm{op}} \le \|\Pi_p\|_{\mathrm{op}}\,\|\Phi_p\|_{\mathrm{op}}
\le \kappa_p.
\end{equation}
\end{lemma}

\begin{proof}
This is the standard submultiplicativity of the operator norm. Since each
$\Pi_p$ is an orthogonal projector, its norm is at most $1$. Therefore
\begin{equation}
\|\Pi_p\Phi_p\|_{\mathrm{op}} \le \|\Pi_p\|_{\mathrm{op}}\,\|\Phi_p\|_{\mathrm{op}}
\le 1\cdot \kappa_p = \kappa_p.
\end{equation}
\end{proof}

\begin{proposition}[Global norm bound for $\Omega$]
Assume the family $\{\Phi_p\}_{p\in\mathbb{P}}$ is uniformly bounded by
$\kappa = \sup_p \kappa_p < \infty$. Then $\Omega$ is bounded on $\mathcal{H}$ and
satisfies
\begin{equation}
\|\Omega\|_{\mathrm{op}} \le \kappa.
\end{equation}
If $\kappa < 1$, then $\Omega$ is a strict contraction.
\end{proposition}

\begin{proof}
Let $x = \bigoplus_p x_p \in \mathcal{H}$. Then
\begin{align}
\|\Omega x\|_{\mathcal{H}}^2
&= \sum_{p\in\mathbb{P}} w_p \|\Pi_p\Phi_p x_p\|_p^2 \\
&\le \sum_{p\in\mathbb{P}} w_p \|\Pi_p\Phi_p\|_{\mathrm{op}}^2 \|x_p\|_p^2 \\
&\le \sum_{p\in\mathbb{P}} w_p \kappa_p^2 \|x_p\|_p^2 \\
&\le \kappa^2 \sum_{p\in\mathbb{P}} w_p \|x_p\|_p^2 \\
&= \kappa^2 \|x\|_{\mathcal{H}}^2.
\end{align}
Taking square roots yields
\begin{equation}
\|\Omega x\|_{\mathcal{H}} \le \kappa \|x\|_{\mathcal{H}}.
\end{equation}
Hence $\|\Omega\|_{\mathrm{op}} \le \kappa$. If $\kappa<1$, the contraction
estimate is strict.
\end{proof}

\begin{corollary}[Banach fixed-point regime]
If $\|\Omega\|_{\mathrm{op}} \le \kappa < 1$, then $\Omega$ has a unique fixed point
$x^\ast \in \mathcal{H}$ and, for every initial state $x_0$, the iterates satisfy
\begin{equation}
\|\Omega^n x_0 - x^\ast\|_{\mathcal{H}} \le \kappa^n \|x_0-x^\ast\|_{\mathcal{H}}.
\end{equation}
\end{corollary}

\begin{proof}
This is the Banach contraction theorem on the complete metric space
$(\mathcal{H},\|\cdot\|_{\mathcal{H}})$.
\end{proof}

\section{Bounds for the PIRTM Update Operator}

At step $t$, the runtime map is
\begin{equation}
\mathcal{U}_t(T) = P F_t K_t(T),
\end{equation}
where $K_t$ is the prime-indexed kernel operator.

\begin{definition}[Kernel bound]
Assume that the kernel operator $K_t : \mathcal{H} \to \mathcal{H}$ is bounded and
satisfies
\begin{equation}
\|K_t\|_{\mathrm{op}} \le \nu_t.
\end{equation}
\end{definition}

\begin{proposition}[PIRTM operator norm bound]
Suppose $P$ is nonexpansive, $\|F_t\|_{\mathrm{op}} \le \mu_t$, and
$\|K_t\|_{\mathrm{op}} \le \nu_t$. Then
\begin{equation}
\|\mathcal{U}_t\|_{\mathrm{op}} \le \mu_t \nu_t.
\end{equation}
In particular, if $\mu_t\nu_t < 1$, the step map is contractive.
\end{proposition}

\begin{proof}
By operator-norm submultiplicativity,
\begin{equation}
\|\mathcal{U}_t\|_{\mathrm{op}} = \|P F_t K_t\|_{\mathrm{op}}
\le \|P\|_{\mathrm{op}}\,\|F_t\|_{\mathrm{op}}\,\|K_t\|_{\mathrm{op}}.
\end{equation}
Since $P$ is nonexpansive, $\|P\|_{\mathrm{op}}\le 1$. Therefore
\begin{equation}
\|\mathcal{U}_t\|_{\mathrm{op}} \le 1\cdot \mu_t \nu_t = \mu_t\nu_t.
\end{equation}
If $\mu_t\nu_t<1$, the map is a strict contraction.
\end{proof}

\begin{corollary}[Uniform runtime contractivity]
If there exists $\rho<1$ such that
\begin{equation}
\mu_t\nu_t \le \rho < 1 \qquad \text{for all } t,
\end{equation}
then every step map $\mathcal{U}_t$ is uniformly contractive and every finite
composition satisfies
\begin{equation}
\|\mathcal{U}_{t+n-1}\circ \cdots \circ \mathcal{U}_t\|_{\mathrm{op}} \le \rho^n.
\end{equation}
\end{corollary}

\begin{proof}
Each factor has norm at most $\rho$, and repeated application of
submultiplicativity yields the result.
\end{proof}

\section{ACE Budget Conservation}

\begin{definition}[ACE observable]
For a state $x\in\mathcal{H}$ define the ACE observable
\begin{equation}
\mathsf{A}(x) = \sum_{p\in\operatorname{supp}(x)} \alpha(p)\,\vartheta_p(x_p).
\end{equation}
\end{definition}

\begin{proposition}[ACE projection is budget preserving]
Assume $P$ is projection onto the feasible region
\begin{equation}
\mathcal{A} = \{x : \mathsf{A}(x) \le B\}.
\end{equation}
Then, for every state $x$,
\begin{equation}
\mathsf{A}(P x) \le B.
\end{equation}
Moreover, if $x\in\mathcal{A}$ already, then $P x = x$.
\end{proposition}

\begin{proof}
The first statement follows from the definition of $P$ as projection into the
feasible set. The second follows from idempotence of projection onto a closed
constraint set: points already feasible are fixed by the projector.
\end{proof}

\begin{proposition}[One-step ACE monotonicity]
Let $T_{t+1}=P F_t K_t(T_t)$. Then
\begin{equation}
\mathsf{A}(T_{t+1}) \le B.
\end{equation}
If, in addition, $T_t$ is feasible and $F_tK_t(T_t)$ does not exceed the current
budget mass before projection, then
\begin{equation}
\mathsf{A}(T_{t+1}) \le \mathsf{A}(T_t).
\end{equation}
\end{proposition}

\begin{proof}
The first inequality is immediate from the previous proposition with
$x = F_tK_t(T_t)$. The second uses the fact that a nonexpansive budget projector
cannot increase the ACE observable when the pre-projected point is already at or
below the current feasible mass. In the concrete radial-scaling implementation,
$P$ rescales by a factor in $[0,1]$, hence every positive homogeneous observable
is non-increasing.
\end{proof}

\begin{remark}
The code-level implementation computes a before/after budget mass and rescales by
\begin{equation}
\lambda = \frac{\text{total before}}{\text{total after} + \varepsilon_0},
\qquad 0<\lambda\le 1,
\end{equation}
whenever the post-kernel mass exceeds budget. This makes ACE monotonicity
explicit rather than merely abstract.
\end{remark}

\section{PWCFL Weighted Contraction Bound}

The weighted average contraction statistic is defined by
\begin{equation}
\mathrm{WAC}_N(t)=
\frac{\sum_{i=t-N}^{t} e^{-\lambda(t-i)} |\delta_i|}
{\sum_{i=t-N}^{t} e^{-\lambda(t-i)}},
\qquad \lambda > 0.
\end{equation}

\begin{proposition}[Convex-hull bound for WAC]
For any finite window,
\begin{equation}
\min_{t-N\le i\le t} |\delta_i| \le \mathrm{WAC}_N(t)
\le \max_{t-N\le i\le t} |\delta_i|.
\end{equation}
\end{proposition}

\begin{proof}
The weights
\begin{equation}
\omega_i = \frac{e^{-\lambda(t-i)}}{\sum_{j=t-N}^{t} e^{-\lambda(t-j)}}
\end{equation}
are strictly positive and sum to one. Therefore $\mathrm{WAC}_N(t)$ is a convex
combination of the values $|\delta_i|$, which must lie between their minimum and
maximum.
\end{proof}

\begin{corollary}[Plateau and divergence tests]
If $|\delta_i|\le \eta$ for all $i$ in the window, then
\begin{equation}
\mathrm{WAC}_N(t)\le \eta.
\end{equation}
If instead $|\delta_i|\ge 1+\varepsilon$ for all $i$ in the window, then
\begin{equation}
\mathrm{WAC}_N(t)\ge 1+\varepsilon.
\end{equation}
\end{corollary}

\begin{proof}
Apply the convex-hull bound with the corresponding uniform lower or upper bound.
\end{proof}

\begin{proposition}[Monotone-window halt justification]
Suppose the sequence $|\delta_{t-N}|,\dots,|\delta_t|$ is nondecreasing and
$|\delta_{t-N}|>1$. Then
\begin{equation}
\mathrm{WAC}_N(t) > 1.
\end{equation}
Hence the PWCFL halt condition is justified whenever the upward trend criterion
coincides with all windowed contractions already exceeding unity.
\end{proposition}

\begin{proof}
By hypothesis every term in the weighted average exceeds $1$, so the weighted
average itself exceeds $1$.
\end{proof}

\section{PWEH Path-Binding and Perturbation Separation}

Let the PWEH chain be defined recursively by
\begin{equation}
H_0 = \mathrm{H}(s_0\|m_0),
\qquad
H_{k+1} = \mathrm{H}(H_k\|s_{k+1}\|m_{k+1}\|\pi(k+1)).
\end{equation}

\begin{proposition}[Prefix sensitivity]
If two executions agree through step $r-1$ and differ at step $r$ in either the
state payload $s_r$, metadata $m_r$, or prime label $\pi(r)$, then the resulting
hash-chain states satisfy
\begin{equation}
H_r \neq H_r'
\end{equation}
except with collision probability of the underlying hash primitive.
\end{proposition}

\begin{proof}
At stage $r$, the chained input string to $\mathrm{H}$ differs between the two
executions. By collision resistance and second-preimage resistance of the hash,
these distinct inputs yield distinct outputs except with negligible probability.
\end{proof}

\begin{corollary}[Full path binding]
If two executions differ at any finite stage, then their terminal hashes differ:
\begin{equation}
H_n \neq H_n'
\end{equation}
for every terminal stage $n$ after the first disagreement, except with negligible
probability.
\end{corollary}

\begin{proof}
By the previous proposition the first divergence forces distinct chained states.
Thereafter all later hashes incorporate distinct previous values, so the chains
remain separated unless a collision occurs.
\end{proof}

\begin{claimbox}
This is the formal reason that NAFL path equivalence in the main article may
require both terminal-state equality and PWEH-chain equality: equality of output
state alone does not guarantee equality of the prime-indexed causal history.
\end{claimbox}

\section{Non-Associative Path Dependence in NAFL}

\begin{definition}[Prime-ordered action]
Let $g_p$ and $g_q$ be prime-indexed operators on a tensor manifold $\mathcal{M}$.
The system is non-associative and order-sensitive if there exist a state
$X\in\mathcal{M}$ and primes $p\neq q$ such that
\begin{equation}
g_p(g_q(X)) \neq g_q(g_p(X)).
\end{equation}
\end{definition}

\begin{proposition}[Distinct ordered paths may yield distinct manifold points]
If the prime-indexed operator family contains at least one noncommuting pair,
then there exist two paths with the same multiset of prime labels but different
orderings that produce different terminal states.
\end{proposition}

\begin{proof}
Take primes $p\neq q$ and a state $X$ such that
$g_p(g_q(X)) \neq g_q(g_p(X))$. Then the two paths $(p,q)$ and $(q,p)$ contain
the same prime multiset but produce different outputs.
\end{proof}

\begin{corollary}[Path history as state]
In an NAFL regime with noncommuting prime actions, the effective cognitive state
cannot be represented solely by the terminal tensor. It must include ordered path
history or an equivalent path-binding certificate.
\end{corollary}

\begin{proof}
If two different paths can produce different semantics and possibly even the same
terminal tensor under coarse observation, then the terminal tensor alone is not a
complete state descriptor. The ordered path or its authenticated digest is needed.
\end{proof}

\section{Spectral Safety Margin and Runtime Guarding}

The runtime implementation enforces a contractivity gate of the form
\begin{equation}
\mathrm{gain}(T_t) = \frac{\rho(T_{t+1})}{\rho(T_t)+\varepsilon_0}
< \operatorname{op\_norm\_T} + \epsilon,
\end{equation}
where $\rho(\cdot)$ denotes a chosen spectral radius surrogate or operator norm.

\begin{proposition}[Sufficient runtime safety margin]
Assume that for a given backend the spectral surrogate obeys
\begin{equation}
\rho(F_tK_t x) \le \mu_t\nu_t\,\rho(x)
\end{equation}
for all nonzero $x$, and that projection $P$ is nonexpansive in the same
surrogate. Then
\begin{equation}
\mathrm{gain}(T_t) \le \mu_t\nu_t + O(\varepsilon_0).
\end{equation}
Hence any configuration with
\begin{equation}
\mu_t\nu_t \le \operatorname{op\_norm\_T} + \epsilon - \delta
\end{equation}
for some fixed $\delta>0$ satisfies the runtime guard for sufficiently small
regularisation constant $\varepsilon_0$.
\end{proposition}

\begin{proof}
By assumption and nonexpansiveness of $P$,
\begin{equation}
\rho(T_{t+1}) = \rho(PF_tK_tT_t) \le \rho(F_tK_tT_t)
\le \mu_t\nu_t\,\rho(T_t).
\end{equation}
Therefore
\begin{equation}
\mathrm{gain}(T_t) = \frac{\rho(T_{t+1})}{\rho(T_t)+\varepsilon_0}
\le \mu_t\nu_t\frac{\rho(T_t)}{\rho(T_t)+\varepsilon_0}
\le \mu_t\nu_t.
\end{equation}
If one tracks finite-precision perturbations or surrogate mismatch, the estimate
is stable up to an $O(\varepsilon_0)$ correction. The claimed margin follows.
\end{proof}

\begin{invariantbox}
This proposition formalises the role of the protected fields
\texttt{epsilon} and \texttt{op\_norm\_T}. They are not generic hyperparameters;
they define an admissible safety envelope for stepwise spectral gain.
\end{invariantbox}

\section{Summary of Explicit Bounds}

For convenience, the principal bounds proved above are collected here.

\begin{table}[h!]
\centering
\caption{Core operator and runtime bounds}
\begin{tabular}{p{5cm}p{6.5cm}}
\toprule
\textbf{Object} & \textbf{Bound} \\
\midrule
Sector map $\Pi_p\Phi_p$ & $\|\Pi_p\Phi_p\|_{\mathrm{op}} \le \kappa_p$ \\
Core operator $\Omega$ & $\|\Omega\|_{\mathrm{op}} \le \sup_p \kappa_p$ \\
PIRTM step $\mathcal{U}_t = PF_tK_t$ & $\|\mathcal{U}_t\|_{\mathrm{op}} \le \mu_t\nu_t$ \\
$n$-step runtime composition & $\|\mathcal{U}_{t+n-1}\circ\cdots\circ\mathcal{U}_t\| \le \rho^n$ if $\mu_t\nu_t\le \rho<1$ \\
ACE projection $P$ & $\|P\|_{\mathrm{op}} \le 1$ \\
WAC statistic & $\min |\delta_i| \le \mathrm{WAC}_N \le \max |\delta_i|$ \\
Runtime spectral gain & $\mathrm{gain}(T_t) \le \mu_t\nu_t + O(\varepsilon_0)$ \\
\bottomrule
\end{tabular}
\end{table}

\section{Formalisation Targets}

The arguments in this appendix are sufficient for a mathematically precise
article appendix, but they also identify exact formalisation targets for the
Lean4 proof pipeline:
\begin{enumerate}
  \item sectorwise boundedness and uniform contraction of $\Omega$,
  \item nonexpansiveness of the ACE projector,
  \item compositional contraction of the PIRTM step map,
  \item correctness of runtime guard inequalities,
  \item path-separation properties of the PWEH chain,
  \item necessity of history retention in NAFL under noncommuting prime actions.
\end{enumerate}

These targets are small enough to be staged independently while still covering
most of the mathematical claims needed for audit-grade verification.

\nocite{*}
\bibliographystyle{unsrt}  
\bibliography{references}  


\end{document}
